\section{Preliminaries}
The technical foundation of this thesis comes from recent progress in transformer-based language modeling, multilingual NLP, retrieval systems, and RAG-based question answering. The Transformer architecture replaced recurrent sequence processing with self-attention, enabling scalable representation learning for natural language \cite{vaswani2017attention}. Modern large language models extend this idea by training on large text corpora and generating fluent responses to user prompts \cite{touvron2023llama,dubey2024llama3}.

For Bangla and other low-resource languages, language-specific and multilingual representation learning is important. BanglaBERT demonstrates that pretraining on Bangla corpora can improve downstream Bangla NLP performance \cite{bhattacharjee2021banglabert}. However, Civic.ai does not currently fine-tune a Bangla language model. Instead, it uses multilingual retrieval and local LLM generation because the main thesis goal is evidence-grounded government-service QA, not training a new foundation model.

\section{Transformer and Large Language Models}
Transformer-based LLMs can understand and generate natural language, but their knowledge is stored parametrically and may become outdated or unsupported. In public-service guidance, this is risky because users may rely on the answer to make decisions about documents, fees, correction procedures, or office visits. LLMs can also hallucinate, meaning they may produce statements that sound plausible but are not grounded in evidence. Verification-oriented work such as Chain-of-Verification shows that explicit checking can reduce hallucination in LLM outputs \cite{dhuliawala2023chain}. For Civic.ai, this supports the decision to use retrieval grounding and controlled answer paths for high-risk civic questions.

\section{Retrieval-Augmented Generation}
Retrieval-Augmented Generation combines two components: a retriever and a generator. The retriever searches an external knowledge base, while the generator produces an answer conditioned on retrieved passages \cite{lewis2020rag}. RAG is especially useful for knowledge-intensive tasks because source documents can be updated without retraining the language model.

Recent RAG surveys describe multiple stages of RAG evolution, including naive RAG, advanced RAG, modular RAG, hybrid retrieval, reranking, query transformation, and evaluation \cite{gao2023rag_survey,gupta2024rag_survey}. These studies show that retrieval quality is often the most important factor in answer correctness. If the retriever returns irrelevant or noisy chunks, even a strong LLM may generate a poor answer. This directly supports the thesis priority: accuracy first, retrieval quality second, reasoning third, and fluency last.

A related bilingual Bangla-English RAG study, NextRAG, shows the relevance of preprocessing, hybrid retrieval, and local LLM comparison in a Bangladeshi QA context \cite{nextrag2025}. Civic.ai is inspired by the general direction of such work, but it applies the design to government-service guidance instead of financial question answering.

\section{Dense Retrieval, Sparse Retrieval, and Hybrid Retrieval}
Dense retrieval represents queries and documents as vectors and retrieves content based on semantic similarity. Dense Passage Retrieval showed that learned dense representations can improve open-domain question answering \cite{karpukhin2020dpr}. Sentence-BERT also demonstrated that sentence-level semantic embeddings are useful for similarity and retrieval tasks \cite{reimers2019sbert}. In Civic.ai, dense retrieval is important because citizen questions are often paraphrased. For example, a user may ask about a lost birth certificate in informal wording that does not exactly match the official FAQ title.

Sparse retrieval, such as BM25, ranks documents based on lexical term matching and term weighting \cite{robertson2009bm25}. BM25 remains useful for exact government-service terms, fee names, legal references, and official terminology. Hybrid retrieval combines dense and sparse signals to reduce the weaknesses of either method alone. Reciprocal Rank Fusion is a rank-level fusion method that combines ranked lists without requiring raw dense and sparse scores to be directly comparable \cite{cormack2009rrf}.

\section{Multilingual and Bangla-Centric Considerations}
A Bangladeshi government-service assistant cannot be designed as an English-only system. Users may ask questions in Bangla, English, or code-mixed Bangla-English. Source documents may contain formal Bangla, English service terms, legal expressions, and OCR or encoding issues. This makes multilingual embeddings important. BGE-M3 is selected because it is designed for multilingual, multi-functionality, and multi-granularity retrieval \cite{chen2024bgem3}. This supports cross-lingual and paraphrase-based retrieval better than an English-only embedding model.

The system also separates \texttt{retrieval\_text} from \texttt{answer\_evidence}. Retrieval text may include aliases and normalized wording to improve search. Answer evidence preserves original Bangla source content so that the LLM is grounded on authentic evidence rather than artificial query-expansion text. This distinction is central to making the system both searchable and faithful.

\section{AI in Government and Civic Services}
Research on AI in government highlights opportunities for efficiency, transparency, and citizen engagement, but it also raises concerns about trust, accountability, and data governance \cite{pislaru2024citizen,alhosani2024ai_government,andrews2022trust}. Public-service chatbots must be more careful than casual assistants because wrong answers may directly affect citizens' decisions. Studies on chatbots in public service also show that user trust depends on transparency, reliability, and clear expectation setting \cite{larsen2024public_service_chatbots,lin2023government_chatbots,mehr2017ai_government}.

For this reason, Civic.ai is designed to show source IDs, retrieval scores, and official/source links. The goal is not only to answer the user but also to make the answer inspectable.

\section{Evaluation of RAG Systems}
Evaluation is essential for thesis-level RAG work. General NLP benchmarks such as GLUE and SuperGLUE evaluate language understanding, while XQuAD and TyDi QA evaluate multilingual and cross-lingual question answering \cite{wang2018glue,wang2019superglue,artetxe2019xquad,clark2020tydiqa}. However, a domain-specific RAG system also needs retrieval-specific and generation-specific evaluation.

Information retrieval metrics such as Hit@k/Recall@k, Mean Reciprocal Rank, and nDCG are used to measure whether expected evidence appears in the retrieved results and how highly it is ranked \cite{manning2008ir,voorhees1999trec8qa,jarvelin2002ndcg}. RAGAS proposes automated metrics such as faithfulness, answer relevance, and context relevance for retrieval-augmented systems \cite{es2024ragas}. When reference answers are available, QA-style exact match or F1 can be used \cite{rajpurkar2016squad}. When wording differs but meaning is similar, semantic metrics such as BERTScore are useful \cite{zhang2020bertscore}.

\section{Research Gap}
Existing work shows strong progress in LLMs, RAG, multilingual NLP, and AI for public service. However, there is still a practical gap for Bangla-English government-service RAG systems built from noisy Bangladeshi documents. Many studies focus on clean datasets, English-heavy benchmarks, or broad chatbot behavior. Civic.ai addresses this gap by focusing on Bangla-heavy government-service documents, document-type-aware chunking, retrieval/generation separation, local LLMs, and measurable evaluation of retrieval alternatives.
